\scanpage{222}
La ``condition d'isomorphie'' (ou première
condition d'homogénéité) dit que l'hom
$G\xrightarrow{\delta}Z$ est surjectif, i.e. que l'on a
une suite exacte courte
\diagram{\begin{tikzcd}[ampersand replacement=\&,column sep=large]
{\text{(98)}} \& 1\arrow[r,dash] \& G_0\arrow[r,dash] \& G\arrow[r,"\delta"] \& Z\arrow[r,dash] \& 1 \& {.}\\
\& \& \& Z\arrow[u,hook]\arrow[ur,swap,"{\lambda\mathrel{\unsure{\mapstochar\relbar|\mapsto}}\lambda^2}"] \& \& \&
\end{tikzcd}}

\deleted{J'aimerais maintenant, dans le contexte}

\deleted{Je suppose maintenant que la structure
de graphe \deleted{\unsure{envis.}} $(S,A_0)$, correspondant à un
$i_0\in\pi$, est triangulaire et même associée
à une structure de $2$-polyèdre \inserted{connexe}, et que
$G_0$ soit transitif sur l'ensemble des
repères. Je vais faire bon poids d'hypothèses,
quitte à raffiner par la suite, en supposant
le $2$-polyèdre $P_0$ orientable, soit $G_0^+$ le sous-groupe
d'indice $2$ de $G_0$ qui conserve l'orientation
-- on \unsure{soit|voit} que c'est le même pour $P_0$ et
pour les $P_i$ ($i\in\pi$). A vrai dire, ce qu'on
va utiliser, c'est la donnée d'un sous-groupe}
\[
 \deleted{G_0^+\subset G_0} \tag*{\deleted{(7)}}
\]
\deleted{qui ait les propriétés}
\[
 \left\{\begin{aligned}
 \deleted{\text{a) } G_0^+ \text{ d'indice }2\text{ dans }G_0,\text{ et invariant dans }G}\\
 \deleted{\text{b) } G_0^+\text{ est \uline{simplement} \uline{transitif} sur l'ens }\overline A_0}\\
 \deleted{\text{des arêtes orientées de }P_0}
 \end{aligned}\right. \tag*{\deleted{(8)}}
\]
% archive sheet number: 221

